\documentclass[../main]{subfiles}
\ifSubfilesClassLoaded{\setcounter{chapter}{10}}{}
\begin{document}

\chapter{Level Design dan Game Mechanics}

\begin{subcpmk}
  \item Sub-CPMK 9.1: Merancang level dengan prinsip level design dan mengimplementasikan spawn system
  \item Sub-CPMK 9.2: Mengimplementasikan game state management, progression system, dan object pooling
\end{subcpmk}

\section{Level Design dan Spawn System}

\begin{learningoutcomes}
\item Memahami prinsip level design
\item Mengimplementasikan spawn system (asteroid, enemy)
\item Mengatur wave dan difficulty
\end{learningoutcomes}

\subsection{Level Design}

Flow: tutorial $\rightarrow$ easy $\rightarrow$ hard. Pacing: tension dan release. Space Survivor: asteroid spawn dari atas, terus meningkat. Spawn point, spawn rate, spawn variety.

\subsection{Spawn System}

Spawner script: InvokeRepeating atau coroutine. Instantiate(prefab, position, rotation). Random position di area spawn. Object Pooling untuk optimasi (reuse daripada Instantiate/Destroy).

\section{Game Mechanics dan Progression}

\begin{learningoutcomes}
\item Merancang core game loop
\item Mengimplementasikan score dan progression
\item Menerapkan game design patterns
\end{learningoutcomes}

\subsection{Core Loop}

Space Survivor: Move $\rightarrow$ Shoot $\rightarrow$ Destroy Asteroid $\rightarrow$ Score $\rightarrow$ Survive. Win condition: survive X detik. Lose condition: hit by asteroid.

\subsection{Progression}

Score system. High score (PlayerPrefs). Power-up (cepat, multi-shot). Difficulty scaling: spawn rate naik over time. Game state: Menu, Playing, Paused, GameOver.


\begin{aktivitas}
  \item Implementasikan AsteroidSpawner dengan coroutine
  \item Tambahkan score system dan high score
  \item Implementasikan difficulty scaling
\end{aktivitas}

\begin{latihan}
  \item Jelaskan prinsip level design untuk game arcade!
  \item Apa keuntungan Object Pooling?
  \item Bagaimana implementasi wave system?
\end{latihan}

\begin{asesmen}
\textbf{Praktik}: Spawn system lengkap dengan progression.
\textbf{Rubrik}: Lihat Lampiran A
\end{asesmen}

\begin{checklist}
  \item Saya dapat mengimplementasikan spawn system
  \item Saya memahami core game loop
  \item Saya dapat menggunakan PlayerPrefs
\end{checklist}

\begin{rangkuman}
Level design: flow, pacing. Spawn: Instantiate, coroutine/InvokeRepeating. Object Pooling untuk optimasi. Score, high score, difficulty scaling. Game state management.
\end{rangkuman}

\ifSubfilesClassLoaded{
  \renewcommand{\bibname}{Daftar Pustaka}
  \bibliographystyle{plain}
  \bibliography{../references}
}{}
\end{document}
